\documentclass[11pt,a4paper]{ltjsarticle}
% Common preamble; compile the four documents with LuaLaTeX.
\usepackage{fix-cm}
\usepackage[margin=25mm]{geometry}
\usepackage{luatexja-fontspec}
\setmainjfont{HaranoAjiMincho-Regular.otf}[BoldFont=HaranoAjiMincho-Bold.otf]
\setsansjfont{HaranoAjiGothic-Medium.otf}[BoldFont=HaranoAjiGothic-Bold.otf]
\usepackage{amsmath,amssymb,amsthm,mathtools,booktabs,enumitem}
\usepackage[unicode,hidelinks]{hyperref}
\setlength{\emergencystretch}{2em}
\setlist{nosep}
\newtheorem{theorem}{定理}[section]
\newtheorem{proposition}[theorem]{命題}
\newtheorem{lemma}[theorem]{補題}
\newtheorem{corollary}[theorem]{系}
\theoremstyle{definition}
\newtheorem{definition}[theorem]{定義}
\newtheorem{remark}[theorem]{注意}
\newtheorem{example}[theorem]{例}
\newtheorem{problem}[theorem]{問題}
\newcommand{\C}{\mathbb C}
\newcommand{\Z}{\mathbb Z}
\newcommand{\Q}{\mathbb Q}
\newcommand{\Cat}{\mathcal A}
\newcommand{\DD}{\mathcal D}
\newcommand{\OO}{\mathcal O}
\newcommand{\TT}{\mathcal T}
\newcommand{\PP}{\mathcal P}
\newcommand{\SSer}{\mathbb S}
\newcommand{\Db}{D^{\mathrm b}}
\newcommand{\Hom}{\operatorname{Hom}}
\newcommand{\End}{\operatorname{End}}
\newcommand{\Ext}{\operatorname{Ext}}
\newcommand{\RHom}{\operatorname{RHom}}
\newcommand{\qgr}{\operatorname{qgr}}
\newcommand{\Gr}{\operatorname{Gr}}
\newcommand{\QGr}{\operatorname{QGr}}
\newcommand{\Tor}{\operatorname{Tor}}
\newcommand{\coh}{\operatorname{coh}}
\newcommand{\thick}{\operatorname{thick}}
\newcommand{\gldim}{\operatorname{gldim}}
\newcommand{\GKdim}{\operatorname{GKdim}}
\newcommand{\rk}{\operatorname{rk}}
\newcommand{\id}{\operatorname{id}}
\newcommand{\Hilb}{\operatorname{Hilb}}
\newcommand{\versiondate}{2026年11月17日・v14}

\renewcommand{\versiondate}{2026年11月18日・v18}
\hypersetup{pdftitle={圏論的Fano指数：研究ノート}}
\title{圏論的Fano指数：研究ノート}
\author{}
\date{\versiondate}
\begin{document}
\maketitle
このノートは現在有効な数値的Serre根と非正規化exact指数の補助情報を整理する。
証明本体は\texttt{papers/normalized\_fano\_index.tex}、次数9認識に関する別の補助ノートは
\texttt{uedakazushi/nc-dp9}に置く。
\tableofcontents

\section{圏論的Fano指数の現行補助資料}
数値格子の半直積定理，周期Euler欠損，三次元の合同条件は
\texttt{normalized\_fano\_index.tex}の定理
\texttt{thm:twisted-lattice-root}，\texttt{cor:periodic-deficit}，
\texttt{thm:threefold-arithmetic}に全文と証明を置く．
ここでは再検証しやすい座標と，圏論的実現に残る障害を記す．

次数6 del PezzoのNSの基底$(H,E_1,E_2,E_3)$では交点行列は
$Q=\operatorname{diag}(1,-1,-1,-1)$，標準類は$K=(-3,1,1,1)^t$である．
六本のtoric境界曲線を巡回させる引戻し作用は
\[
 F=\begin{pmatrix}
 2&1&1&1\\-1&-1&0&-1\\-1&-1&-1&0\\-1&0&-1&-1
 \end{pmatrix},\qquad
 F^6=I,\quad F^tQF=Q,\quad FK=K.
\]
$H_0=E_1$とすると$N_6(F)H_0=-K$．
この根のEuler軌道は$k=0,1,\ldots,6$でそれぞれ
$(1,1,2,3,4,5,7)$，非正周期欠損は
$(0,-7/12,-1/3,-1/4,-1/3,-7/12,0)$である．
これらの行列・全次数の式は\texttt{verify\_normalized\_index.py}で正確に検算できる．

\emph{次の障害：} 数値の$F$と$H$が存在しても，それらを同時に実現する
heart-exact自己同値は一般には構成できない．
幾何的格子で$H^1(C_r,\operatorname{NS})$が計算できても，
実際のPicard群や双加群の根の障害を表しているとは限らない．
非可換射影平面のrank三では数値正規化が自動であるが，
一般六公理の圏がrank三のgeometric surface-like擬格子を持つこと自体が未証明である．

三次元の数値候補はSerre双対性，exceptional構造対象，漸近Euler三次性，
軌道正規化を仮定する．Loginov--Orlovの2026年二論文が扱うEuler degreeと本研究の$d$の
同定は未検証であり，数値条件は非可換$M_{72}$の根の実在証明でも反例でもない．

\section{非正規化exact指数の幾何学的試験例}
指数原稿\texttt{normalized\_fano\_index.tex}の
定理\texttt{thm:delpezzo-exact-index-table}に，滑らかなdel Pezzo曲面の
非正規化exact指数の次数別最大値と全証明をまとめた．この研究ノートに証明を重複させない．
次数2のGeiser involution，次数3のFermat cubic，次数4の四点blowup，
次数5の$\mathfrak S_5$，次数6のtoric hexagonは，今後の非可換根の実現性比較のbenchmarkである．
次数4の具体例は$f(u,v)=(v,u^{-1})$の軌道
$(0,t),(t,\infty),(\infty,t^{-1}),(t^{-1},0)$（$t^4\ne1$）をblowupする．
四点を通る$(1,1)$曲線の有無を判定する行列式は非零の座標倍を除き$(t^4-1)/t^2$である．
次数5では$A_4$型根格子のCoxeter元（5-cycle）と例外曲線の五本の軌道の和が$-K$になる．
これらの有限次元の因子類計算は\texttt{verify\_del\_pezzo\_roots.py}で検算する．
このスクリプトは個々の幾何学的自己同型の存在証明を代替しない．
次数4では一般の$X$の自己同型群が$(\Z/2)^4$であり，exact四乗根は存在しない一方，特殊な四乗根による数値格子作用は同じ格子上へ移せる．よって整数数値根がexact自己同値に持ち上がらない実例になる（指数原稿\texttt{prop:degree-four-lifting-gap}）．
この区別は数値根の実現問題に本質的なので，今後の障害理論のbenchmarkとする．
既知幾何との文献上の新規性および非可換例への拡張は未確認．


\end{document}
